\section{Limitations and Threats to Validity}
\label{sec:limitations}

\subsection{Label Leakage via Attack-Specific Parameters}
\label{sec:leakage}
The most significant limitation is the presence of features that directly encode the attack type and intensity. As demonstrated by SHAP (Section \ref{sec:discussion}) and ablation (Table \ref{tab:ablation}), the top features for tree models are:
\begin{itemize}
    \item \texttt{swap\_intensity}, \texttt{n\_swaps} (swap\_injection only)
    \item \texttt{gate\_ins\_intensity}, \texttt{n\_insertions} (gate\_insertion only)
    \item \texttt{dp\_error\_rate} (depolarizing\_noise only)
    \item \texttt{mt\_flip\_prob} (measurement\_tampering only)
    \item \texttt{pl\_detuning\_mag} (pulse\_level\_attack only)
    \item \texttt{tr\_t1}, \texttt{tr\_t2}, \texttt{tr\_gate\_time} (thermal\_relaxation only)
    \item \texttt{has\_noise\_model} (binary: clean vs. any attack)
\end{itemize}

These features allow the classifier to learn simple decision rules (e.g., ``if \texttt{swap\_intensity} $> 0$ then swap\_injection'') rather than detecting attack signatures from quantum measurement statistics. The \textbf{14-19\% accuracy drop} (17-22\% F1 drop) when removing the attack-specific parameter group quantifies the leakage. A controlled experiment removing only \texttt{has\_noise\_model} showed minimal impact on tree models (RF: 95.45\% $\to$ 95.50\%, XGB: 95.75\% $\to$ 95.50\%), confirming that the attack-specific parameters (not the binary indicator) are the primary leakage vector.

\subsection{Different Threat Model}
QAD simulates an adversary who \textbf{modifies the quantum circuit} before or during execution (gate tampering, noise injection). The reference paper addresses an adversary who \textbf{intercepts the quantum channel} between Alice and Bob (photon interception, basis manipulation, detector attacks). These are fundamentally different:
\begin{itemize}
    \item QAD: White-box circuit access, parametric attacks
    \item BB84: Black-box channel access, physical-layer attacks
\end{itemize}

\subsection{No Temporal Dimension}
QAD provides single aggregate statistics per circuit execution (1024 shots). The paper's innovation is temporal-QBER: a sequence of QBER values computed over windows of $W=50$ sifted bits, capturing error dynamics within a session. Our \texttt{features.py} implements the 63 temporal features but there is no temporal data to apply them to.

\subsection{No Basis Reconciliation}
QAD measures all qubits in the computational basis after circuit execution. BB84 uses random basis selection (rectilinear/diagonal) with basis sifting, discarding $\sim$50\% of bits. This basis-dependent structure enables the paper's basis-dependent feature category (10 features).

\subsection{No Photon Loss / Detection Efficiency Model}
QAD uses fixed 1024 shots per circuit. Real QKD systems have:
\begin{itemize}
    \item Channel loss (exponential with distance)
    \item Detector efficiency ($<100\%$)
    \item Dark counts, afterpulsing
    \item Variable sifted key length per session
\end{itemize}
The paper's channel-interaction features (16 features) model these effects.

\subsection{Artificially Balanced Classes}
QAD balanced dataset has exactly 100,000 samples per attack (300,000 clean). Real attack frequencies are unknown and likely highly imbalanced. The unbalanced variant exists but was not used.

\subsection{Simulator-Specific Noise Models}
QAD uses Qiskit AerSimulator noise models (depolarizing, amplitude damping, thermal relaxation). These may not reflect:
\begin{itemize}
    \item Hardware-specific noise correlations
    \item Non-Markovian noise
    \item Crosstalk, calibration drift
    \item Environmental fluctuations
\end{itemize}

\subsection{Single Experimental Run}
Unlike the paper's 10 independent seeds, we report results from a single train/test split (random\_state=42) and single model training per algorithm. Variance estimates come only from 5-fold CV.

\subsection{No Baseline Comparison}
The fixed 11\% QBER threshold and CUSUM baselines require temporal QBER sequences, which QAD does not provide. Our implementation includes the framework but cannot compute these baselines for QAD.

\subsection{Reproducibility Scope}
While our pipeline is reproducible (fixed seeds, locked dependencies, saved artifacts), exact reproduction of the reference paper's results is \textbf{not possible} without:
\begin{enumerate}
    \item The paper's exact 63 feature formulas (only categories given)
    \item The paper's attack simulation parameters (not fully specified in 6 pages)
    \item The paper's BB84 simulator implementation
\end{enumerate}